\documentclass[11pt,a4paper,landscape]{ctexart}

\usepackage[top=1.5cm,bottom=1.5cm,left=1.6cm,right=1.6cm]{geometry}
\usepackage{amsmath,amssymb}
\usepackage{booktabs}
\usepackage{array}
\usepackage{float}
\usepackage{xcolor}
\usepackage{enumitem}
\usepackage{hyperref}
\hypersetup{hidelinks}

\pagestyle{empty}

\newcommand{\slidetitle}[2]{%
    \noindent{\Large\bfseries #1}\hfill{\normalsize\color{gray!70!black} #2}\par
    \vspace{0.2em}
    \noindent\textcolor{gray!50}{\rule{\linewidth}{0.6pt}}\par
    \vspace{0.6em}
}

\begin{document}

% ================================
% 封面
% ================================
\begin{center}
\vspace*{4.5cm}

{\fontsize{28}{34}\selectfont\bfseries
掩码感知的神经控制微分方程}

\vspace{1.2em}

{\Large\itshape
用于缺失数据条件下电池 SOH 估计的一个方案}

\vspace{3cm}

\noindent\textcolor{gray!50}{\rule{\linewidth}{0.5pt}}\\[0.3em]
{\normalsize\bfseries 核心问题：}\\[0.2em]
{\normalsize 现有 NCDE 方法在插值后丢弃缺失掩码，}\\[0.1em]
{\normalsize 长连续缺失窗口内的插值平滑段被当作真实观测。}\\[0.6em]
\noindent\textcolor{gray!50}{\rule{\linewidth}{0.5pt}}\\[0.3em]
{\normalsize\bfseries 直接思路：}\\[0.2em]
{\normalsize 将缺失掩码构造为置信度路径，与观测路径联合输入 NCDE，}\\[0.1em]
{\normalsize 同时对隐藏状态施加缺失时长依赖的衰减。}\\[0.3em]
\noindent\textcolor{gray!50}{\rule{\linewidth}{0.5pt}}

\end{center}

\newpage

% ================================
% Page 2: 现有方法做了什么 vs 掩码带来了什么
% ================================
\slidetitle{现有 NCDE 方法的共同操作}{}

\vspace{0.8em}

\begin{minipage}[t]{0.56\textwidth}
\normalsize
\textbf{给定数据：} \(\{(t_i, x_i, m_i)\}_{i=1}^N\)

\(x_i\)：测量值（电压、电流、温度等）\\
\(m_i \in \{0,1\}\)：缺失指示（0=缺失，1=观测）

\vspace{0.8em}
\textbf{现有 NCDE（BatteryCDE, PA-NCDE 等）：}
\begin{enumerate}[leftmargin=1.5em, itemsep=0.6em]
    \item 对 \(x_i\) 做样条或线性插值 → 得到连续路径 \(X(t)\)
    \item \(X(t)\) 输入 NCDE 做隐藏状态演化
    \item 输出 SOH 估计
\end{enumerate}

\vspace{0.6em}
\textbf{信息流中被丢弃的部分：}
\[
\{m_i\} \;\longrightarrow\; \boxed{\text{丢弃}}
\]

\vspace{0.4em}
模型不知道：\\
-- 哪些点是真实观测的\\
-- 插值段距离最近观测点有多远\\
-- 缺失是否与工况相关

\end{minipage}\hfill
\begin{minipage}[t]{0.42\textwidth}
\vspace{0em}
\normalsize
\textbf{缺失掩码携带的信息：}

\vspace{0.6em}
\begin{enumerate}[leftmargin=1.2em, itemsep=0.8em]
    \item \textbf{点的真伪}\\
    区分"传感器真实读数"和"插值猜测值"。

    \item \textbf{距最近观测的时长}\\
    \(\Delta t(t) = t - t_{\text{last\_obs}}\)： \\
    时长越长，\(X(t)\) 的偏差风险越大。

    \item \textbf{缺失与工况的关联}（MNAR）\\
    大电流/高温下丢包率更高时，\\
    缺失本身暗示系统处于严苛状态。
\end{enumerate}
\end{minipage}

\newpage

% ================================
% Page 3: 两个修改点
% ================================
\slidetitle{两个修改点}{双通道路径 + 状态衰减}

\vspace{0.6em}

\normalsize

\textbf{修改点 1：掩码增强的双通道路径}

\vspace{0.4em}
现有做法：输入路径 \(X(t) \in \mathbb{R}^d\)（仅插值后的观测值）

修改后：
\[
\text{输入路径：}\quad Z(t) = [\,X(t)\;\; C(t)\,] \in \mathbb{R}^{d+1}
\]

其中 \(C(t)\) 由掩码构造，缺失区间内 \(C(t) \in [0,1)\)：
\[
C(t) = \exp(-\gamma \cdot \Delta t_{\text{last\_obs}}(t)) \qquad \text{或对 } m_i \text{ 做线性插值}
\]

\vspace{1.5em}
\textbf{修改点 2：掩码引导的状态遗忘}

\vspace{0.4em}
标准 NCDE 演化：
\[
dh(t) = f_\theta(h(t))\,dZ(t)
\]

加入遗忘项（当 \(C(t) \approx 0\) 时激活）：
\[
dh(t) = f_\theta(h(t))\,dZ(t) \;-\; \alpha \cdot (1-C(t)) \odot h(t)\,dt
\]

长缺失区间内隐藏状态指数衰减至默认状态，重新获得观测后从保守状态重启。

\vspace{1.5em}
\textbf{属性：}

\vspace{0.4em}
\begin{itemize}[leftmargin=2em, itemsep=0.4em]
    \item 两项修改是\textbf{插件式的}：不改变 NCDE 基础架构，不改变数据预处理
    \item 当缺失率 $=0$ 时 \(C(t)\equiv 1\)，模型退化为标准 NCDE
    \item 参数增量极少（仅 \(C(t)\) 可学习的衰减系数），计算开销可忽略
\end{itemize}

\newpage

% ================================
% Page 4: 预期行为
% ================================
\slidetitle{不同缺失模式下的预期行为}{对比标准 NCDE（无掩码）}

\vspace{0.6em}

\begin{table}[H]
\centering
\normalsize
\renewcommand{\arraystretch}{1.4}
\begin{tabular}{@{}p{4.5cm} p{10.5cm} p{10.5cm}@{}}
\toprule
\textbf{缺失模式} & \textbf{标准 NCDE} & \textbf{MA-NCDE 预期} \\
\midrule
\textbf{随机缺失} &
各点等概率，插值段短。 &
\textbf{差异不大}。\(C(t)\approx 1\) 大多时刻，不造成副作用。 \\
\midrule
\textbf{连续块缺失} &
插值填充整个缺失窗口，模型误认为平滑观测。 &
衰减机制持续降低隐藏状态。恢复观测后从低状态重启，\textbf{预期优势最大}。 \\
\midrule
\textbf{MNAR} &
不感知丢包上下文。 &
\(C(t)<1\) 隐含指示数据不可靠，\textbf{可能有可测优势}。\\
\bottomrule
\end{tabular}
\end{table}

\vspace{0.4em}
\begin{center}
\normalsize
随机缺失（无信息）\quad $\approx$\quad
块缺失（有结构）\quad $<$\quad
MNAR（含信号）
\end{center}

\newpage

% ================================
% Page 5: 范围与已知未知
% ================================
\slidetitle{范围与已知未知}{}

\vspace{0.8em}

\normalsize

\textbf{这个方案做了什么：}

\vspace{0.3em}
\begin{itemize}[leftmargin=2em, itemsep=0.4em]
    \item 将缺失掩码从"被丢弃的信息"变为连续置信度路径
    \item 在连续时间框架内对不可靠区间做状态衰减
    \item 两项修改是对标准 NCDE 的插件式补充
\end{itemize}

\vspace{1em}
\textbf{当前没有讨论的问题：}

\vspace{0.3em}
\begin{itemize}[leftmargin=2em, itemsep=0.4em]
    \item \(C(t)\) 的具体构造形式（线性插值 vs 指数衰减 vs 可学习参数）\hfill --- 待实验
    \item 衰减系数 \(\alpha\) 的敏感性和最优值\hfill --- 待实验
    \item 是否加入辅助的掩码重构损失\hfill --- 可选，暂未计入
    \item 跨数据集的泛化表现\hfill --- 可能在 NASA 数据上先验证
\end{itemize}

\newpage

% ================================
% 封底
% ================================
\begin{center}
\vspace*{7cm}
{\Huge\bfseries 方案结束}
\end{center}

\end{document}
